\section{Physics reference}
\label{sec:physics}

This appendix collects the formulae implemented in the
\code{cls\_estimations/} backend and used throughout the workflow. It
is intentionally compact; for derivations and detailed discussion,
see SATLAS2~\cite{satlas2} and the Handbook of Nuclear Physics
chapter on nuclear spins and moments~\cite{cls_handbook}.

\subsection{Doppler shift}

For an ion accelerated through a potential difference $V$ with
charge state $q$ and rest mass $m$, the relativistic Lorentz factor
is
\begin{equation}
    \gamma = 1 + \frac{qV}{m c^{2}}, \qquad
    \beta = \sqrt{1 - \gamma^{-2}}.
\end{equation}
The lab-frame laser frequency $\nu_\text{lab}$ relates to the
frequency seen by the ion in its rest frame, $\nu_\text{rest}$, by
\begin{equation}
    \nu_\text{rest} = \gamma\,\nu_\text{lab}\,(1 - \beta\,\hat{n} \cdot \hat{k})
                    = \gamma\,\nu_\text{lab}\,(1 \mp \beta),
    \label{eq:doppler}
\end{equation}
where the upper (lower) sign in the right-most expression
corresponds to anti-collinear (collinear) geometry. Equation
\eqref{eq:doppler} is what \denis{} inverts when the user specifies a
target lab-frame voltage offset \app{dV} from $V_0$ for a chosen
isotope --- the laser frequency $\nu_\text{lab}$ is back-solved so
that the desired rest-frame centroid sits at $V = V_0 - \mathrm{dV}$.

\subsection{Hyperfine structure}

For a level with electronic angular momentum $J$ and nuclear spin
$I$, the total angular momentum $F = I + J$ takes values
$|I-J|, |I-J|+1, \dots, I+J$. The hyperfine energy of the substate
$F$ is
\begin{equation}
    \Delta E_F = \frac{A}{2}\, C
    + B\,\frac{\frac{3}{2}\,C(C+1) - 2I(I+1)\,J(J+1)}
                {2 I (2I-1)\,J(2J-1)},
    \label{eq:hfs}
\end{equation}
with
\begin{equation}
    C = F(F+1) - I(I+1) - J(J+1).
\end{equation}
$A$ is the magnetic-dipole coupling constant and $B$ the
electric-quadrupole coupling constant; the second term in
\eqref{eq:hfs} vanishes when $I < 1$ or $J < 1$. \denis{} uses
\eqref{eq:hfs} for both lower and upper electronic states; the
observed transition frequencies are differences between the two.
Selection rules are $\Delta F = 0, \pm 1$ with $F = 0 \to F = 0$
forbidden.

\subsection{Schmidt magnetic moments}

For a single odd nucleon in an orbital with quantum numbers
$\ell, j$ the Schmidt prediction is
\begin{equation}
    \mu = j\,g_j\,\mu_N, \qquad
    g_j = \begin{cases}
        g_\ell + (g_s - g_\ell)/(2\ell+1), & j = \ell + 1/2 \\
        g_\ell - (g_s - g_\ell)\,\dfrac{j+1}{2j+1}, & j = \ell - 1/2,
    \end{cases}
\end{equation}
with bare-nucleon $g$-factors $g_\ell^{(p)} = 1$, $g_\ell^{(n)} = 0$,
$g_s^{(p)} = +5.586$, $g_s^{(n)} = -3.826$. The optional $g_s$
quenching factor in the calculator multiplies the spin $g$-factor to
account for in-medium effects. The two-particle case couples the
two single-particle moments through the Lande formula at a chosen
total spin $I$.

\subsection{Fit objectives}

The \app{Fitter} block exposes three statistics:

\begin{description}
    \item[Chi-square.]
        $\chi^2 = \sum_i \left(\frac{y_i - m_i(\theta)}{\sigma_i}\right)^2$,
        with $\sigma_i$ taken from the \app{yerr} model
        ($\sqrt{y_i}$, $\sqrt{y_i+1}$, or model-based). Suitable for
        bins with at least $\sim 10$ counts.
    \item[Poisson LLH.]
        $-2 \ln L = 2 \sum_i \big[m_i(\theta) - y_i + y_i \ln(y_i/m_i(\theta))\big]$
        for $y_i > 0$. Use this when the per-bin counts are low and
        the Gaussian approximation breaks down.
    \item[Gaussian LLH.]
        Same kernel as $\chi^2$ but with $\sigma_i$ allowed to depend
        on the model. Useful when the noise model is heteroskedastic.
\end{description}

\app{Scale covariance} on the \app{Fitter} block multiplies the
post-fit covariance matrix by the reduced $\chi^2$, recovering the
standard errors when the input $\sigma_i$ are known to be
mis-scaled.

\subsection{Statistical methods and systematics}

\paragraph{Fit statistics.} $\chi^2$ with $\sqrt{y_i+1}$ errors is
biased at low counts; the Poisson likelihood is not~\cite{baker1984}.
Its numerical error bars come out $\sqrt{2}$ too large unless the
\app{Fitter} block's correction is ticked.

\paragraph{Overdispersion.} For Poisson counts the variance of a step's
counts over the scans equals their mean. \app{Run quality} reports
variance/mean $D$ from each step's repeats, with the significance of
$D-1$ from the Poisson null, and splits a PMT's excess into the part all
PMTs share (the ions per bunch) and its own (afterpulsing), which the
dead-time gate removes.

\paragraph{Weighted means.} Runs are averaged with weights $1/\sigma^2$
and the error scaled by $\max(1, \sqrt{\chi^2_\nu})$, the Birge
ratio~\cite{birge1932}.

\paragraph{Validate fit.} A parametric bootstrap: spectra drawn from the
best fit (at the run's measured $D$ when \app{Run quality} ran) and
refitted with the same settings. Each parameter's pull
$(\hat\theta-\theta)/\sigma$ should have mean 0 and width 1, and the
$\pm1\sigma$ / $\pm2\sigma$ intervals should cover the truth 68\,\% /
95\,\% of the time.

\paragraph{The systematic budget.} Each source in the Systematics tab's
\app{Sources} table is a $1\sigma$ width on one perturbation of the
apparatus (cooler offset or divider, scan divider, scan-axis offset or
gain, laser offset), with a kind (measured, declared, unknown) and a
correlation class (\code{run}, \code{campaign}, \code{region},
\code{table}). \app{Propagate sources} refits the whole chain with each
source at $\pm1\sigma$. A source that moves every run at once enters a
weighted mean as the weighted mean of the run responses, sign kept; a
per-run source as their weighted quadrature sum. On request, Monte Carlo
draws the shared sources together from $N(0,\sigma)$; per-run sources
still enter from their $\pm1\sigma$ responses. The method follows
A.~Raggio's review of the analysis (November 2026); the drift correction
follows~\cite{vandenborne2025}.

\subsection{Nuclear moments and charge radii}

\paragraph{Ratio method.} Against a reference isotope on the same line,
\begin{equation}
    \mu = \mu_\mathrm{ref}\,\frac{I}{I_\mathrm{ref}}\,\frac{A}{A_\mathrm{ref}},
    \qquad
    Q = Q_\mathrm{ref}\,\frac{B}{B_\mathrm{ref}},
\end{equation}
from each level where measured and their mean, with the fit's
correlation of $A_l$ and $A_u$. The hyperfine anomaly is neglected
unless a percentage is given, which enters as an error.

\paragraph{Charge radii.} With
$1/\mu_{AA'} = (m_{A'}-m_A)/(m_A m_{A'})$ from the ions' masses in the
charge state measured,
\begin{equation}
    \delta\nu^{AA'} = F\,\delta\langle r^2\rangle^{AA'} + K/\mu_{AA'},
    \qquad
    \delta\langle r^2\rangle = (\delta\nu - K/\mu)/F .
\end{equation}
$F$ and $K$ are given, calibrated by a York fit~\cite{york2004} of the
modified shifts $\mu\,\delta\nu$ against $\mu\,\delta\langle r^2\rangle$
for pairs of known radii, or taken from a second line's
shifts~\cite{king1984}. Shifts against one reference share its error,
and the tab keeps the full covariance. An isomer pair has $1/\mu = 0$:
its shift is $F\,\delta\langle r^2\rangle$ alone, and it is left off the
King plot.
